%%%%%%%%%%%%%%%%%%%%%%%%%%%%%%%%%%%%%%%%%
% Medium Length Professional CV
% LaTeX Template
% Version 2.0 (8/5/13)
%
% This template has been downloaded from:
% http://www.LaTeXTemplates.com
%
% Original author:
% Rishi Shah 
%
% Important note:
% This template requires the resume.cls file to be in the same directory as the
% .tex file. The resume.cls file provides the resume style used for structuring the
% document.
%
%%%%%%%%%%%%%%%%%%%%%%%%%%%%%%%%%%%%%%%%%

%----------------------------------------------------------------------------------------
% PACKAGES AND OTHER DOCUMENT CONFIGURATIONS
%----------------------------------------------------------------------------------------

\documentclass{resume} % Use the custom resume.cls style

\usepackage[left=0.75in,top=0.6in,right=0.75in,bottom=0.6in]{geometry} % Document margins
\newcommand{\tab}[1]{\hspace{.2667\textwidth}\rlap{#1}}
\newcommand{\itab}[1]{\hspace{0em}\rlap{#1}}
\name{Banjou Tahara} % Your name

% クラスファイルの仕様上、最初に入力した address が一番下に、最後に入力した address が一番上に印刷されます。
\address{tahara.banjou.rwq@ecs.osaka-u.ac.jp} % Your email
\address{5-36 Shikinai-cho, Izumiotsu-shi, Osaka 595-0035, Japan} % Your address


\begin{document}

%----------------------------------------------------------------------------------------
% EDUCATION SECTION
%----------------------------------------------------------------------------------------
\begin{rSection}{Education}

  \begin{rSubsection}{The University of Osaka, Japan}{April 2026 -- Present}{Master's Program}{Department of Systems Innovation, Graduate School of Engineering Science}
    \item Enrolled in the Graduate School of Engineering Science.
  \end{rSubsection}

  \begin{rSubsection}{The University of Osaka, Japan}{April 2022 -- March 2026}{Bachelor of Engineering}{Department of Systems Science, School of Engineering Science}
    \item Graduated First in the Department (Valedictorian).
  \end{rSubsection}

\end{rSection}

%----------------------------------------------------------------------------------------
% CAREER OBJECTIVE SECTION
%----------------------------------------------------------------------------------------
\begin{rSection}{Career Objective}
  \item To advance next-generation spatial computing by shifting the paradigm from rigid physical limitations and isolated VR/AR environments to a ubiquitous Real-Cyber co-existence. Leveraging a rigorous multidisciplinary background, I aim to synergize computational imaging, the psychophysics of material appearance, and HCI. My ultimate goal is to dissolve the boundaries between digital and physical domains, bringing the boundless flexibility of Cyber-Worlds into the real physical world to elevate human quality of life (QoL).
\end{rSection}

%----------------------------------------------------------------------------------------
%   PROJECTS SECTION
%----------------------------------------------------------------------------------------
\begin{rSection}{Projects}

  \begin{rSubsection}{Spatial Overwriting via Volumetric Scanning Projection Mapping}{2025 - Present}{}{}
    \item Developed a novel Projection Mapping (PM) framework to realistically manipulate physical material appearance (color, gloss, and texture) without permanent surface modification. Addressed the fundamental limitation of flat-looking projections by introducing a volume-scanning approach to reproduce "mesoscopic" 3D micro-structures (e.g., fine bumps and roughness).
    \smallskip

    \item \textbf{Proposed Method \& Results:} Designed a volumetric scanning projection system that applies rapid, microscopic 2D orbital vibrations to the projection target. By synchronizing these movements with the high-speed switching of normal-mapped depth layer images, the system successfully enables the perception of textures that bring out 3D bump geometries within the temporal integration window of the human visual system.

    \item \textbf{Future Vision:} To address current technical challenges, immediate future work focuses on refining hardware precision to eliminate motion blur caused by the induced vibrations. Following this optimization, I plan to conduct user studies and rigorous psychophysical evaluations to investigate the relationship between the 3D displacement geometry data of textures and human material perception.
  \end{rSubsection}

\end{rSection}



%----------------------------------------------------------------------------------------
%   PUBLICATIONS AND PRESENTATION LIST
%----------------------------------------------------------------------------------------

\begin{rSection}{Publications and Presentations}

  % --- 日本語・査読なし論文 ---
  \begin{rSubsection}{Domestic Conference / Non-Peer-Reviewed}{}{}{}
    \item 1. Banjou Tahara, Daisuke Iwai, "Geometric Relief Representation via High-Speed Volumetric Projection synchronized with 2D Micro-Vibration", The 70th Annual Conference of the Institute of Systems, Control and Information Engineers (SCI'26), Japan, pp.772-778, March.24- March.26, 2026.
  \end{rSubsection}

\end{rSection}



%----------------------------------------------------------------------------------------
% ACADEMIC ACHIEVEMENTS SECTION
%----------------------------------------------------------------------------------------
\newpage
\begin{rSection}{Academic Awards}

  \begin{rSubsection}{Kusumoto Award (Valedictorian / Top Graduating Student)}{March 2026}{}{}
    \item Awarded by The University of Osaka to the top graduating student in the Department of Systems Science, School of Engineering Science, for outstanding academic excellence.
  \end{rSubsection}

  \begin{rSubsection}{SICE Outstanding Student Award}{March 2026}{}{}
    \item Conferred by The Society of Instrument and Control Engineers (SICE) to recognize students who demonstrate exceptional academic performance and exemplary character.
  \end{rSubsection}

\end{rSection}



%----------------------------------------------------------------------------------------
% TECHNICAL STRENGTHS SECTION
%----------------------------------------------------------------------------------------


\begin{rSection}{Technical Strengths}

  \begin{list}{}{
      \setlength{\leftmargin}{1.5em}
      \setlength{\topsep}{0pt}
    }
    \item \begin{tabular}{ @{} >{\bfseries}l @{\hspace{4ex}} l }
            Core Disciplines   & Computer Vision, Computer Graphics, Optics, Visual Psychology, \\
                               & Signal Processing, Measurement \& Control Engineering,         \\
                               & Human-Computer Interaction, Applied Physics,                   \\
                               & Electronic Engineering, System Optimization.                   \\
            \noalign{\smallskip}
            Graphics \& Vision & OpenCV, OpenGL, Physically Based Rendering (PBR),              \\
                               &
            %  Ray Tracing, 
            Projectior-Camera Systems, Projection-Mapping,                                      \\
                               & Material Appearance Science.                                   \\
            \noalign{\smallskip}
            AI \& Robotics     & PyTorch, Multimodal World Models for Robot Learning,           \\
                               & Generative AI, Image Tokenizer Training.                       \\
            \noalign{\smallskip}
            Sensor Fusion      & Multi-Sensor Integration (LiDAR, IMU, GNSS)                    \\
                               & for Autonomous Driving Development.                            \\
          \end{tabular}
  \end{list}

\end{rSection}



%----------------------------------------------------------------------------------------
% WORK EXPERIENCE SECTION
%----------------------------------------------------------------------------------------
\begin{rSection}{Intern Experience}

  \begin{rSubsection}{Matsuo Institute, Inc.}{May 2025 - April 2026}{Research Engineer (Intern)}{Tokyo, Japan}
    \item Engaged in cutting-edge research and development on \textbf{Multimodal World Models} for Robot Learning.
    \item Developed and trained custom image tokenizers using PyTorch to build robust neural network baselines for robotic decision-making.
    \item Implemented large-scale machine learning pipelines to synthesize multimodal sensor data, bridging generative AI with physical robotic control.
  \end{rSubsection}

  \begin{rSubsection}{KG Motors Co., Ltd.}{June 2024 - December 2025}{Autonomous Driving Software Engineer (Intern)}{Hiroshima, Japan}
    \item Developed an end-to-end \textbf{Autonomous Driving System} deployed on a physical vehicle through the integration and synchronization of heterogeneous sensors (LiDAR, IMU, and GNSS).
    \item Implemented a high-fidelity \textbf{Digital Twin} environment synchronized with the physical system to simulate and validate vehicle kinematics, sensor behaviors, and path-planning logic.
  \end{rSubsection}

\end{rSection}

%----------------------------------------------------------------------------------------
% EXTRA-CURRICULAR SECTION
%----------------------------------------------------------------------------------------
\newpage
\begin{rSection}{Extra-Curricular \& Competitions}

  \begin{rSubsection}{Autonomous Driving AI Challenge 2025 (JSAE)}{July 2025 - October 2025}{}{}
    \item Competed in an autonomous driving lap time racing competition organized by the Society of Automotive Engineers of Japan (JSAE).
    \item Achieved \textbf{4th place in the Qualifying Round} and \textbf{7th place in the Final Round} out of 249 participating teams.
    \item Developed and optimized path-planning algorithms and vehicle control logic to minimize lap times while ensuring safe obstacle avoidance.
  \end{rSubsection}

  % \begin{rSubsection}{Yosakoi Dance Team "MADANI"}{April 2022 - September 2024}{}{}
  %   \item Performed as an active dancer and core member in The University of Osaka's representative traditional-modern fusion festival dance team.
  %   \item Advanced to the \textbf{finals at the Koiya Festival} (September 2024), one of the largest regional Yosakoi dance competitions in the Kansai area.
  %   \item Contributed to team synchronization, large-scale choreography formations, and intense physical conditioning, fostering strong leadership and collaborative teamwork under high-pressure live environments.
  % \end{rSubsection}
\end{rSection}

%----------------------------------------------------------------------------------------
% PERSONAL TRAITS SECTION
%----------------------------------------------------------------------------------------
\begin{rSection}{Personal Traits}

  \item \textbf{Growth-Oriented Research Mindset (Positive Thinking):} Views technical setbacks and experimental failures not as roadblocks, but as essential data points and catalyst for growth. Applies a resilient, constructive mindset to overcome the inherent challenges of pioneering research.

  \item \textbf{Synergistic Hardware-Software Intuition:} Possesses a rare, hands-on understanding of the synergy between hardware constraints and software optimization. Experienced in bridging the gap between digital algorithms (computer vision, ML, digital twins) and physical mechanics (projection systems, real-vehicle heterogeneous sensors).

\end{rSection}

\end{document}



























% %%%%%%%%%%%%%%%%%%%%%%%%%%%%%%%%%%%%%%%%%
% % Medium Length Professional CV
% % LaTeX Template
% % Version 2.0 (8/5/13)
% %
% % This template has been downloaded from:
% % http://www.LaTeXTemplates.com
% %
% % Original author:
% % Rishi Shah 
% %
% % Important note:
% % This template requires the resume.cls file to be in the same directory as the
% % .tex file. The resume.cls file provides the resume style used for structuring the
% % document.
% %
% %%%%%%%%%%%%%%%%%%%%%%%%%%%%%%%%%%%%%%%%%

% %----------------------------------------------------------------------------------------
% % PACKAGES AND OTHER DOCUMENT CONFIGURATIONS
% %----------------------------------------------------------------------------------------

% \documentclass{resume} % Use the custom resume.cls style

% \usepackage[left=0.75in,top=0.6in,right=0.75in,bottom=0.6in]{geometry} % Document margins
% \newcommand{\tab}[1]{\hspace{.2667\textwidth}\rlap{#1}}
% \newcommand{\itab}[1]{\hspace{0em}\rlap{#1}}
% \name{Banjou Tahara} % Your name

% % クラスファイルの仕様上、最初に入力した address が一番下に、最後に入力した address が一番上に印刷されます。
% % 意図した順序（上が住所、下が連絡先）にするため、記述順を入れ替えています。
% \address{tahara.banjou.rwq@ecs.osaka-u.ac.jp} % Your phone number and email
% \address{5-36 Shikinai-cho, Izumiotsu-shi, Osaka 595-0035, Japan} % Your address


% \begin{document}
% %----------------------------------------------------------------------------------------
% % EDUCATION SECTION
% %----------------------------------------------------------------------------------------

% % \begin{rSection}{Education}

% %   {\bf Veermata Jijabai Technological Institute, Mumbai} \hfill {\em August 2018 - Present}
% %   \\ Master in Technology
% %   \\ Department of Structural Engineering\\

% %   {\bf Maharashtra Institute of Technology, Pune} \hfill {\em July 2013 - June 2017}
% %   \\ Bachelor of Engineering, Civil.\hfill { Overall Percentage: 68.14 }

% % \end{rSection}

% \begin{rSection}{Education}

%   {\bf The University of Osaka，Japan} \hfill {\em April 2026 -- Present}
%   \\ Master's Program
%   \\ Department of Systems Innovation， Graduate School of Engineering Science

%     {\bf The University of Osaka，Japan} \hfill {\em April 2022 -- March 2026}
%   \\ Bachelor of Engineering
%   \\ Department of Systems Science，School of Engineering Science
%   \\ Graduated First in the Department

% \end{rSection}

% %----------------------------------------------------------------------------------------
% % CAREER OBJECTIVE SECTION
% %----------------------------------------------------------------------------------------

% % \begin{rSection}{Career Objective}
% %   To work for an organization which provides me the opportunity to improve my skills and knowledge to grow along with the organization objective.
% % \end{rSection}

% \begin{rSection}{Career Objective}
%   To advance the field of spatial computing by seamlessly integrating VR/AR realities and dissolving the flexibility of Cyber-Worlds into Physical-Reality. I aim to leverage Computational Imaging, Science of Material Perception, and HCI to develop stress-free XR interfaces, ultimately enhancing the QoL across diverse industries, from healthcare to digital fabrication.
% \end{rSection}

% % 仮想現実／拡張現実をシームレスに統合し、サイバー世界の柔軟性を物理現実に融合させることで、空間コンピューティング分野の発展に貢献します。計算画像処理、物質外観科学、ヒューマンコンピュータインタラクション（HCI）を活用し、ストレスフリーなクロスリアリティインターフェースを開発することで、医療からデジタルファブリケーションまで、多様な産業における生活の質（QoL）の向上を目指します。

% %----------------------------------------------------------------------------------------
% % PROJECTS AND SEMINARS SECTION
% %----------------------------------------------------------------------------------------


% %----------------------------------------------------------------------------------------
% %   PROJECTS SECTION
% %----------------------------------------------------------------------------------------

% \begin{rSection}{Projects}

%   \begin{rSubsection}{Spatial Overwriting via Volumetric Scanning Projection Mapping}{2025 - Present}{}{}
%     Developed a novel Projection Mapping (PM) framework to realistically manipulate physical material appearance (color, gloss, and texture) without permanent surface modification. Addressed the fundamental limitation of flat-looking projections by introducing a volume-scanning approach to reproduce "mesoscopic" 3D micro-structures (e.g., fine bumps and roughness).
%     \smallskip \\

%     \item \textbf{Proposed Method \& Results:} Designed a volumetric scanning projection system that applies rapid, microscopic 2D orbital vibrations to the projection target. By synchronizing these movements with the high-speed switching of normal-mapped depth layer images, the system successfully enables the perception of textures that bring out 3D bump geometries within the temporal integration window of the human visual system.\\

%     \item \textbf{Future Vision:} To address current technical challenges, immediate future work focuses on refining hardware precision to eliminate motion blur caused by the induced vibrations. Following this optimization, I plan to conduct user studies and rigorous psychophysical evaluations to investigate the relationship between the 3D bump geometry data of textures and human material perception. The ultimate goal is to establish a high-fidelity spatial computing technology that can tangibly elevate human QoL.
%   \end{rSubsection}

% \end{rSection}


% % \begin{rSection}{Projects}
% % \begin{rSubsection}
% %   {\bf Spatial Overwriting via Volumetric Scanning Projection Mapping} \hfill {\em 2025 - Present}
% %   \\ Developed a novel Projection Mapping (PM) framework to realistically manipulate physical material appearance (color, gloss, and texture) without permanent surface modification. Addressed the fundamental limitation of flat-looking projections by introducing a volume-scanning approach to reproduce "mesoscopic" 3D micro-structures (e.g., fine bumps and roughness).

% %   \item \textbf{Proposed Method \& Results:} Designed a volumetric scanning projection system that applies rapid, microscopic 2D orbital vibrations to the projection target. By synchronizing these movements with the high-speed switching of normal-mapped depth layer images, the system successfully enables the perception of textures that bring out 3D bump geometries within the temporal integration window of the human visual system.
% %   \item \textbf{Future Vision:} To address current technical challenges, immediate future work focuses on refining hardware precision to eliminate motion blur caused by the induced vibrations. Following this optimization, I plan to conduct user studies and rigorous psychophysical evaluations to investigate the relationship between the 3D bump geometry data of textures and human material perception. The ultimate goal is to establish a high-fidelity spatial computing technology that can tangibly elevate human QoL.


% % \end{rSection}

% % \begin{rSection}{Projects}
% %   {\bf Dynamic Analysis of Buckling Restrained Braces}
% %   \\The project aims at designing and fabrication of two Buckling Restrained Braces which were analyzed under dynamic loading. As alternative for conventional braces, these BRBs are also beneficial for seismic retro-fitting in RCC and steel structures.\\

% %   {\bf Indirect Model Analysis of Structures}\\
% %   Presented a Seminar on Indirect Model Analysis, explaining the method to compute response of Prototype from the Influence lines obtained from Model. Use of Muller Breslau Principle in Indirect Model Analysis and the Similitude between prototype and model.\\

% %   {\bf Microtunneling}\\
% %   Presented a seminar on Micro Tunneling, explaining its advantages over conventional method of drainage laying systems. Analysis considering direct and indirect cost of micro tunneling was also discussed.

% % \end{rSection}

% %----------------------------------------------------------------------------------------
% % TECHNICAL STRENGTHS SECTION
% %----------------------------------------------------------------------------------------

% %----------------------------------------------------------------------------------------
% % TECHNICAL STRENGTHS SECTION
% %----------------------------------------------------------------------------------------

% \newpage
% \begin{rSection}{Technical Strengths}

%   \begin{tabular}{ @{} >{\bfseries}l @{\hspace{4ex}} l }
%     Core Disciplines   & Computer Vision, Computer Graphics, Optics, Visual Psychology, \\
%                        & Signal Processing, Measurement \& Control Engineering,         \\
%                        & Human-Computer Interaction, Applied Physics,                   \\
%                        & Electronic Engineering, System Optimization                    \\
%     \noalign{\smallskip}
%     Graphics \& Vision & OpenCV, OpenGL, Physically Based Rendering (PBR),              \\
%                        & Ray Tracing, Projection Mapping-Camera Systems,                \\
%                        & Material Appearance Science                                    \\
%     \noalign{\smallskip}
%     AI \& Robotics     & Generative AI, Multimodal World Models for Robot Learning,     \\
%                        & PyTorch, Image Tokenizer Training                              \\
%     \noalign{\smallskip}
%     Sensor Fusion      & Multi-Sensor Integration (LiDAR, IMU, GNSS)                    \\
%                        & for Autonomous Driving Development                             \\
%   \end{tabular}

% \end{rSection}


% %----------------------------------------------------------------------------------------
% % WORK EXPERIENCE SECTION
% %----------------------------------------------------------------------------------------

% \begin{rSection}{Intern Experience}

%   \begin{rSubsection}{Matsuo Institute, Inc.}{May 2025 - April 2026}{Research Engineer (Intern)}{Tokyo, Japan}
%     \item Engaged in cutting-edge research and development on \textbf{Multimodal World Models for Robot Learning}.
%     \item Developed and trained custom image tokenizers using PyTorch to build robust neural network baselines for robotic decision-making.
%     \item Implemented large-scale machine learning pipelines to synthesize multimodal sensor data, bridging generative AI with physical robotic control.
%   \end{rSubsection}

%   \begin{rSubsection}{KG Motors Co., Ltd.}{June 2024 - December 2025}{Autonomous Driving Software Engineer (Intern)}{Hiroshima, Japan}
%     \item Developed an end-to-end \textbf{autonomous driving system deployed on a physical vehicle} through the integration and synchronization of heterogeneous sensors (LiDAR, IMU, and GNSS).
%     \item Implemented a high-fidelity \textbf{digital twin} environment synchronized with the physical system to simulate and validate vehicle kinematics, sensor behaviors, and path-planning logic.
%   \end{rSubsection}

% \end{rSection}


% %----------------------------------------------------------------------------------------
% % ACADEMIC ACHIEVEMENTS SECTION
% %----------------------------------------------------------------------------------------


% \begin{rSection}{Academic Achievements}

%   \textbf{Kusumoto Award (Valedictorian / Top Graduating Student)} \hfill {\em March 2026}
%   \\ Awarded by Osaka University to the top graduating student in the Department of Systems Science, School of Engineering Science, for outstanding academic excellence. \\

%   \textbf{SICE Outstanding Student Award} \hfill {\em March 2026}
%   \\ Conferred by the Society of Instrument and Control Engineers (SICE) to recognize students who demonstrate exceptional academic performance and exemplary character.

% \end{rSection}

% \newpage

% %----------------------------------------------------------------------------------------
% % EXTRA-CURRICULAR SECTION
% %----------------------------------------------------------------------------------------

% \begin{rSection}{Extra-Curricular \& Competitions}

%   \begin{rSubsection}{Autonomous Driving AI Challenge 2025 (JSAE)}{\em July 2025 - October 2025}{}{}
%     \item Competed in an autonomous driving lap time racing competition organized by the Society of Automotive Engineers of Japan (JSAE).
%     \item Achieved \textbf{4th place in the Qualifying Round} and \textbf{7th place in the Final Round} out of 249 participating teams.
%     \item Developed and optimized path-planning algorithms and vehicle control logic to minimize lap times while ensuring safe obstacle avoidance.
%   \end{rSubsection}

%   \begin{rSubsection}{Yosakoi Dance Team "MADANI" （祭楽人）}{\em April 2022 - September 2024}{}{}
%     \item Performed as an active dancer and core member in one of Osaka University's premier traditional-modern fusion festival dance teams.
%     \item Advanced to the \textbf{Finals at the Koiya Festival} (September 2024), one of the largest regional Yosakoi dance competitions in the Kansai area.
%     \item Contributed to team synchronization, large-scale choreography formations, and intense physical conditioning, fostering strong leadership and collaborative teamwork under high-pressure live environments.
%   \end{rSubsection}

% \end{rSection}

% % \begin{rSection}{Extra-Curricular}
% %   \begin{itemize}
% %     \item Co-Organized ``Nirmitee 2017'' - a National Symposium of Civil Department of MIT, Pune
% %     \item Attended a workshop on Autodesk Revit at IIT Bombay in 2014.
% %     \item Winner of Inter Departmental Football Competition 2015.
% %     \item Member of the Rotaract Club Of Pune Pride from 2014 to 2017.
% %     \item Worked for a start-up company Named OUST as a Regional Marketing Manager
% %   \end{itemize}
% % \end{rSection}

% %----------------------------------------------------------------------------------------
% % PERSONAL TRAITS SECTION
% %----------------------------------------------------------------------------------------

% \begin{rSection}{Personal Traits}
%   \item \textbf{Growth-Oriented Research Mindset (Positive Thinking):} Views technical setbacks and experimental failures not as roadblocks, but as essential data points and catalyst for growth. Applies a resilient, constructive mindset to overcome the inherent challenges of pioneering research.
%   \item \textbf{High-Energy Resilient Collaboration:} Successfully bridges the professional rigor of corporate R\&D teams (AI/robotics) with the high-stakes, high-pressure competitive environment of elite festival dance formations. Proven capability to harness collective team energy and maintain a goal-oriented mindset to achieve concrete results under high pressure.
%   \item \textbf{Synergistic Hardware-Software Intuition:} Possesses a rare, hands-on understanding of the synergy between hardware constraints and software optimization. Experienced in bridging the gap between digital algorithms (computer vision, ML, digital twins) and physical mechanics (projection systems, real-vehicle heterogeneous sensors).
% \end{rSection}

% \end{document}